\section{Measured results}

All figures are from the \dSoc{} on a \dBoard{} with the firmware of commit \texttt{\dFirmwareCommit}, and they
come from the logs saved as evidence. They are not from the ESP32-S3 nor from the u-blox receiver of the final
PCB.

\subsection{Time of \texttt{geofence\_check()} (PoC-01)}

Mean per call of a batch of 40, in three consecutive boots. The three boots agree within a maximum difference
of 0.02~\% and all 40 classifications of each polygon were correct.

\begin{table}[H]
\centering
\small
\begin{tabularx}{\textwidth}{|r r r X|}
\hline
\tablehead{\textbf{Vertices} & \textbf{Mean per call} & \textbf{Margin against 50~ms} & \textbf{Status}} \\ \hline
4 & \dBenchUsFour~$\mu$s & \dBenchMarginFour$\times$ & Met \\ \hline
16 & \dBenchMsSixteen~ms & \dBenchMarginSixteen$\times$ & Met \\ \hline
64 & \dBenchMsSixtyFour~ms & \dBenchMarginSixtyFour$\times$ & Met \\ \hline
256 & \dBenchMsTwoFiftySix~ms & \dBenchMarginTwoFiftySix$\times$ & Met \\ \hline
1024 & \dBenchMsThousand~ms & \dBenchMarginThousand$\times$ & \textbf{Not met} \\ \hline
\end{tabularx}
\caption{Mean time of \texttt{geofence\_check()} on the board.}
\end{table}

The cost is linear: $T(N) \approx \dBenchFitOffsetUs~\mu\text{s} + \dBenchFitSlopeUs~\mu\text{s}\cdot N$, with
differences of at most 1.5~$\mu$s from the means. The 50~ms limit is reached near \dBenchLimitVertices{}
vertices. \textbf{Criterion PoC-01 is met for polygons of up to about 800 vertices and is not met with 1024.}
The fence of a farm, with tens or a few hundred vertices, falls inside. It is a measurement with the kernel
counter and not with an oscilloscope.

\subsection{Cold-start first fix (PoC-07)}

\begin{table}[H]
\centering
\small
\begin{tabularx}{\textwidth}{|c r X|}
\hline
\tablehead{\textbf{Boot} & \textbf{Time to green LED} & \textbf{Ratio to the \dTtffSpecCold~s of the datasheet}} \\ \hline
1 & \dTtffRunA~s & 3.0$\times$ \\ \hline
2 & \dTtffRunB~s & 1.8$\times$ \\ \hline
3 & \dTtffRunC~s & 4.5$\times$ \\ \hline
\end{tabularx}
\caption{Time to the first usable fix after a cold start, measured with a stopwatch.}
\end{table}

The mean is \dTtffMean~s, the median \dTtffMedian~s and the standard deviation \dTtffSd~s, with only three
samples. The \dSoc{} datasheet gives \dTtffSpecCold~s typical for a cold start, under open sky, at 25~°C and
with an external low-noise amplifier with a SAW filter; it is not known whether this board meets that
assumption. The tests were done in a neighborhood street between houses, not under open sky, so the
comparison is indicative. The uncertainty of the measurement is about 2~s. The time measured is to a fix with
a reported accuracy of 50~m or better, not to the very first fix.

With the same cold start, this time from the window of the static test, \textbf{there was no fix in
\dWindowColdMinutes{} minutes}. In that time the receiver's satellite list ranged between
\dWindowColdListedMin{} and \dWindowColdListedMax, with four or more only \dWindowColdFourPct~\% of the time,
except for a peak of \dWindowColdPeak{} in the first seconds of the search. That list counts satellites that
appear in the receiver's response, not satellites with a usable signal nor used in the fix. A possible
explanation, which was not verified, is that without assistance data the receiver has to decode the orbit of
each satellite from its own signal, and that requires a stronger and more continuous signal than the partly
obstructed sky of that window.

\subsection{GNSS position in the static test (PoC-02 and PoC-03)}

The test was done from a first-floor window facing south, open, with the board practically outside, with the
normal firmware and the stored satellite data. The board stayed still for the \dStaticRecordingMinutes{}
minutes of the recording. The sequence was:

\begin{itemize}
  \item At \dStaticFirstFixS~s after the reset the first fix arrived, with a reported accuracy of
        \dStaticFirstFixAcc~m, which the application discarded for exceeding 50~m. At \dStaticFenceS~s an
        acceptable fix created the fence and the application moved to \texttt{MONITORING}.
  \item \textbf{First run:} until \dStaticRunOneEndS~s, \dStaticRunOneFixes{} accepted fixes at 1~Hz
        (\dStaticRunOneMinutes~min), with a single gap of \dStaticMaxGapS~s. The reported accuracy converged
        from \dStaticAccStart~m to about \dStaticAccConverged~m in under a minute. At the end it degraded
        gradually, from \dStaticTailAccFrom{} to \dStaticTailAccTo~m, the satellite list dropped to 3 or 4, and
        at \dStaticLostFirstS~s the application moved to \texttt{FIX\_LOST}.
  \item \textbf{No fix for \dStaticGapMinutes{} minutes.} The receiver kept searching with the board still.
  \item \textbf{Second run:} at \dStaticRunTwoStartS~s the fix came back, with an initial reported accuracy of
        \dStaticRunTwoFirstAcc~m, and there were \dStaticRunTwoFixes{} accepted fixes until
        \dStaticRunTwoEndS~s (\dStaticRunTwoMinutes~min). At \dStaticLostSecondS~s it degraded again and the
        application went back to \texttt{FIX\_LOST}.
\end{itemize}

That the fix is lost and comes back with the board still indicates that what changes is the signal that
arrives, not the board. The cause was not determined: the PoC firmware does not log the signal strength of
each satellite. One hypothesis, not verified, is that the satellites, as they move across the sky, pass behind
the trees, the cables or the edge of the window.

\textbf{PoC-03.} The interval between fixes within each run was 1.01~s: the receiver delivers 1~Hz in
continuous mode, and therefore \textbf{does not exceed 1~Hz}, as SR-10 asks. That requirement targets the
receiver of the PCB.

\textbf{PoC-02.} The first \dStaticSkipS~s of fixes of each run were discarded, because of the receiver's
convergence, and the scatter of the positions around their own mean was computed. \dStaticBothUsed{} fixes
remained, about 9.5 minutes of data, below the 15 minutes planned.

\begin{table}[H]
\centering
\small
\begin{tabularx}{\textwidth}{|X r r r|}
\hline
\tablehead{\textbf{Statistic} & \textbf{Run 1} & \textbf{Run 2} & \textbf{Both runs}} \\ \hline
Fixes used & \dStaticRunOneUsed & \dStaticRunTwoUsed & \dStaticBothUsed \\ \hline
CEP50 (m) & \dStaticRunOneCepFifty & \dStaticRunTwoCepFifty & \dStaticBothCepFifty \\ \hline
CEP95 (m) & \dStaticRunOneCepNinetyFive & \dStaticRunTwoCepNinetyFive & \dStaticBothCepNinetyFive \\ \hline
DRMS (m) & \dStaticRunOneDrms & \dStaticRunTwoDrms & \dStaticBothDrms \\ \hline
Reported accuracy, median (m) & \dStaticRunOneAcc & \dStaticRunTwoAcc & \dStaticBothAcc \\ \hline
\end{tabularx}
\caption{Scatter of the position in the static test, after discarding the first 60~s of each run.}
\end{table}

With only the fixes whose reported accuracy was \dStableLimit~m or better (\dStableUsed{} fixes from the two
runs), the CEP50 is \dStableCepFifty~m, the CEP95 \dStableCepNinetyFive~m and the DRMS \dStableDrms~m;
filtering by reported accuracy does not improve the scatter.

Against the SRD criterion of \dReqCep~m CEP, the measured CEP50 is \dStaticRatioTwo{} times larger in the best
run, \dStaticRatioOne{} times in the worst and \dStaticRatioBoth{} times with both together. Against the
\dSoc{} datasheet, which gives a CEP50 of \dNordicCepFifty~m in continuous tracking under open sky, it does not
match either, but the conditions are not the same. \textbf{Criterion PoC-02 is not met under these
conditions.} This does not allow saying what it would give under open sky.

Three remarks about the numbers:

\begin{itemize}
  \item \textbf{The scatter of the two runs together is larger than that of each one.} The means of the two
        runs are \dStaticRunSeparation~m apart. That is, the error has a slow part, which changes with the
        visible constellation, and a short test underestimates it. This suggests that the value of an isolated
        run, such as the 3.5~m of the second one, is not representative.
  \item \textbf{The tail is heavy:} the CEP95 is about 2.7 times the CEP50 in each run, that is, most fixes are
        close to the mean but there are large deviations.
  \item \textbf{The accuracy reported by the receiver} (1 sigma, median of about 11 to 13~m) is larger than the
        scatter observed within a run (DRMS of 5 to 8~m). They do not contradict each other: the scatter around
        the mean does not see a constant error that affects all fixes of the run equally.
\end{itemize}

\subsection{LED alert and demo mode (PoC-04, PoC-05 and PoC-11)}

The demo mode log shows \texttt{SEARCHING} $\to$ \texttt{DEMO}, the creation of the 60~m fence around the
fictitious point, 42 steps of the route and \texttt{DEMO} $\to$ \texttt{SEARCHING} when pressed again. Some
steps, as they came out in the log:

\begin{table}[H]
\centering
\small
\begin{tabularx}{\textwidth}{|r c c X|}
\hline
\tablehead{\textbf{Distance to the boundary} & \textbf{Inside} & \textbf{Color (R, G, B)} & \textbf{LED}} \\ \hline
30.00~m & yes & 0, 255, 0 & Steady green \\ \hline
24.00~m & yes & 0, 255, 0 & Steady green \\ \hline
18.00~m & yes & 26, 229, 0 & Starts mixing toward red \\ \hline
12.00~m & yes & 102, 153, 0 & Mix \\ \hline
6.00~m & yes & 179, 76, 0 & Much redder \\ \hline
0.00~m & no & 255, 0, 0 & Blinking red (on the exact boundary the half-open rule gives outside) \\ \hline
6.00~m & no & 255, 0, 0 & Blinking red, outside \\ \hline
4.00~m & yes & 204, 51, 0 & On the way back, inside \\ \hline
10.00~m & yes & 128, 127, 0 & Mix \\ \hline
22.00~m & yes & 0, 255, 0 & Steady green \\ \hline
\end{tabularx}
\caption{Steps of the demo mode (\texttt{evidence/logs/demo\_mode\_session.log}).}
\end{table}

The sequence matches the design: steady green at 20~m or more from the boundary, a mix toward red when
approaching, and blinking red outside, up to 10~m beyond the boundary, and back. The person doing the test
pressed the button, saw the board enter the demo mode and, when pressing again, saw the LED go back to
blinking blue. The demo mode worked without a GNSS signal. \textbf{PoC-04, PoC-05 and PoC-11 are met} in the
PoC.

\subsection{Fail-safe behavior (PoC-06)}

\begin{itemize}
  \item In the static test, the log shows \texttt{MONITORING} $\to$ \texttt{FIX\_LOST} at \dStaticLostFirstS~s
        and again at \dStaticLostSecondS~s, in each case about ten seconds after the last accepted fix, and
        \texttt{FIX\_LOST} $\to$ \texttt{MONITORING} when the fix came back.
  \item The person doing the tests observed the LED blinking magenta when the fix was lost, different from the
        blinking red of being outside the fence.
  \item The PC tests check that a non-finite or negative distance, a non-positive \texttt{safe\_m} and any
        unknown state give the fail-safe result, and that \texttt{FAULT} is only left through the demo mode.
\end{itemize}

\subsection{Memory footprint (PoC-08)}

The normal firmware takes \dFlashBytes~B of flash (\dFlashPercent~\% of the \dFlashRegion~B available to the
application) and \dRamBytes~B of RAM (\dRamPercent~\% of \dRamRegion~B). The build with the benchmark takes
\dFlashBench~B and the one that deletes the satellite data, \dFlashCold~B; both use the same RAM. The fence
engine takes \dEngineBytesCmThree~B compiled for Cortex-M33 and \dEngineBytesEsp~B for ESP32-S3.

\subsection{Algorithm correctness and portability (PoC-09 and PoC-10)}

The engine passes its \dHostTests{} tests (\dHostChecks{} checks) and \dStressCases{} randomized cases with a
fixed seed, without failures, also under memory-error detectors. The worst error of the engine's distance
against haversine was \dStressWorstErrorPct~\%, ten times smaller than the documented 0.2~\%. \textbf{PoC-09 is
met.}

The five pure C modules and the engine compile unchanged, with strict warnings as errors, for the PC, the
Cortex-M33 and the ESP32-S3. \textbf{PoC-10 is met as compilation}; nothing was run on the ESP32-S3.
